\documentclass[../../main]{subfiles}

\begin{document}

\section{Minimal coupling}
\label{sec:minimal_coupling}
Throughout this section, we will closely follow the arguments and derivations of David Tong's book. We will refer to other sources explicitly, if statements do not originate from there. 

The equation of motion for a nonrelativistic particle of mass $m$ and charge $q$ moving in an electromagnetic field is given by the Lorentz force law
\begin{equation}
    m \ddot{\bm{x}} = q(\bm{E} + \dot{\bm{x}} \times \bm{B}).
\end{equation}
Here and below, we often omit the arguments $(\bm x, t)$ of the fields. We use the same charge convention as Tong, so $q=-e$ for an electron, where $e>0$. As god intended, we set \(c = \hbar = 1\) thorughout. Lorentz's law can be derived by the Euler-Lagrange equations from the following Lagrangian
\begin{equation}
    L = \frac 12 m  \dot{\bm{x}}^2 + q \dot{\bm{x}} \cdot \bm{A} - q\phi , 
\end{equation}
where $\bm E = -\bm \nabla \phi - \partial_t \bm A$ and $\bm B = \bm \nabla \times \bm A$. The canoncial momentum is give as $\bm{p} = \frac{\partial L}{\partial \dot{\bm{x}}} = m \dot{\bm{x}} + q\bm{A}$. We can obtain the corresponding Hamiltonian by using the Legendre transformation, which finally yields
\begin{equation}
    H = \dot{\bm{x}} \cdot \bm{p} - L = \frac{1}{2m}(\bm{p} - q\bm{A})^2 + q\phi ,
\end{equation}
with \(\bm{p}=-i\hbar\bm{\nabla}=-i\bm{\nabla}\) in the position representation of quantum mechanics. Note, that this Hamiltonian has the known form of kinetic energy plus potential energy with the exception that the momentum is not $\bm{p}$ but $\bm{p}-q\bm{A}$. Replacing $\bm{p} \to \bm{p} - q\bm{A}$ in order to include the effect of an electromagnetic field inside a Hamiltonian is referred to \emph{minimal coupling}. It is not called maximal coupling, since it does not include all possible types of interaction, e.g. the coupling of a particles intrinsic magnetic momentum $\bm \mu$ with the magnetic field via $-\bm \mu \cdot \bm B$. On the other hand, it is the minimal framwork, that ensures that the Schrödinger equation is gauge covariant under the local $U(1)$ gauge transformations (Altland and Simmons, chapter 5.3.4)
\begin{align}
    \bm{A} &\to \bm{A}' = \bm{A}+\bm{\nabla}\chi \notag \\
    \phi   &\to \phi' = \phi-\partial_t\chi \\
    \psi   &\to \psi' = e^{iq\chi} \psi , \notag
\end{align} \label{eq:minimal_coupling_gauge_transformations}
where $\chi=\chi(\bm{x},t)$ is a smooth real function. The name minimal coupling does hence not seem like such a big stretch.

\subsection{Gauge transformations}
\label{subsec:minimal_coupling_gauge_transformations}

Let's check the aforementioned claim of gauge covariance. The primed and unprimed potentials describe the same electric and magnetic fields. Thus, they must yield the same physical predictions. The transformation $\psi' = U_\chi \psi$, with \(U_\chi=e^{iq\chi(\bm{x},t)}\), already preserves the position probability density, $|\psi'(\bm x, t)|^2 = |\psi(\bm x, t)|^2$. To check the dynamics, we show that whenever $\psi$ solves the original Schrödinger equation, $\psi'$ solbes the equation with the gauge transformed fields. We write $H' = H \left[\bm A', \phi' \right]$.

The time dependent Schrödinger equations reads
\begin{align*}
  i\partial_t\psi &= H\psi
    = \left[\frac{1}{2m}\bigl(-i\bm{\nabla}-q\bm{A}\bigr)^2+q\phi\right]\psi .
\end{align*}
Applying the product rule gives
\begin{align*}
  \bm{p}(U_\chi\psi)
    &= (-i)(iq)\bigl(\bm{\nabla}\chi(\bm{r},t)\bigr)U_\chi\psi
       +(\bm{p}\psi)U_\chi \\
    &= U_\chi\bigl(q(\bm{\nabla}\chi)+\bm{p}\bigr)\psi
\end{align*}
and therefore
\begin{align*}
  (\bm{p}-q\bm{A}')\psi'
    &= (\bm{p}-q\bm{A}-q\bm{\nabla}\chi)U_\chi\psi \\
    &= U_\chi(\bm{p}-q\bm{A})\psi ,
\end{align*}
which is valid for all physical, i.e. sufficiently smooth, $\psi$. We can apply this identity a second time by writing $\Pi_k = p_k - qA_k$ and $\Pi'_k = p_k - qA'_k$ where $k$ denotes the Cartesian coordinate. The identity above therefore reads $\Pi'_k (U_\chi \tilde \psi) = U_\chi \Pi_k \tilde \psi$. Taking $\tilde \psi = \Pi_k \psi$ and summing over $k$ reads
\begin{equation}
    (\bm{p}-q\bm{A}')^2\psi'
       =\sum_k \Pi'_k(U_\chi\Pi_k\psi)
       =U_\chi\sum_k \Pi_k^2\psi
       =U_\chi(\bm{p}-q\bm{A})^2\psi. \label{eq:minimal_coupling_momentum_identity}
\end{equation}
We finally obtain
\begin{align*}
    i\partial_t\psi'
       &=U_\chi\bigl(i\partial_t\psi-q(\partial_t\chi)\psi\bigr) \\
       &=U_\chi\left[\frac{1}{2m}(\bm{p}-q\bm{A})^2
                       +q(\phi-\partial_t\chi)\right]\psi \\
       &=\left[\frac{1}{2m}(\bm{p}-q\bm{A}')^2+q\phi'\right]\psi'
        =H'\psi'.
\end{align*}
where we have applied the sum rule again in the first line, the original Schrödinger equation in the second one and \cref{eq:minimal_coupling_momentum_identity} in the last line.

Thus a gauge transformation maps each solution to a solution of the transformed equation. This is what it means for the Schrödinger equation to be \emph{gauge covariant}. Rewriting the second and third line gives
\begin{align*}
    H' U_\chi \psi = U_\chi (H -q \, \partial_t \chi)\psi 
\end{align*}
for all $\psi$. The Hamiltonian hence transforms as
\begin{equation*}
    H' = U_\chi H U_\chi^\dagger - q\, \partial_t \chi .
\end{equation*}
For a time-independent gauge transformation the two Hamiltonians are unitarily equivalent and therefore have the same spectrum.

\subsection{The Peierls substitution}

Now how do creation and annihilation operators in second quantization transform under a gauge transformation? We will answer this question for the discrete case, when one is looking at models that are defined on a lattice. This section is merely meant to develop some intuition for the origins of the Peierls phase and we refer to the papers \emph{Electromagnetic coupling in tight-binding models for strongly correlated light and matter} by Li et al. and \emph{Gauge fixing for strongly correlated electrons coupled to quantum light} by Dmytruk and Schiró for more rigorous approaches to include light matter coupling in lattice models.

Take a general single-particle state on the lattice in position representation
\begin{equation*}
    \ket{\psi} 
    = \sum_j \braket{\bm r_j}{\psi} \ket{\bm r_j} 
    = \sum_j \psi(\bm r_j) c^\dagger_{\bm r_j} \ket{0}
\end{equation*}
and apply a gauge transformation, which results in
\begin{equation*}
    \ket{\psi'} 
    = \sum_j \psi'(\bm r_j) c^\dagger_{\bm r_j} \ket{0}
    = \sum_j \psi(\bm r_j) U_\chi c^\dagger_{\bm r_j} \ket{0}.
\end{equation*}
We can see, that transforming the creation operators as
\begin{equation}
    c^\dagger_j \to c'^\dagger_j = U_\chi c^\dagger_j = e^{iq\chi(\bm r_j, t)} c^\dagger_j  
\end{equation}
is equivalent to transforming the wave function itself. This is the starting point of the short derivation done in “Tight-binding models in a magnetic field: Peierls substitution,” by Akhmerov, Sau and van Heck. I will subsequently follow their line of reasoning.


  

\paragraph{References}
\begin{itemize}
    \item Tong, David. "Applications of quantum mechanics." University of Cambridge Part II Mathematical Tripos (2017). (Chapter 6.1)
    \item Altland and Simons. Condensed matter field theory. Cambridge university press, 2010. (Chapter 5.3.4)
    \item Bartelmann, et al. Theoretische Physik 2| Elektrodynamik. Berlin: Springer Spektrum, 2018. (Chapter: Lagrange- und Hamilton-Formalismus in der Elektrodynamik)
    \item Bartelmann, et al. Theoretische Physik 3| Quantenmechanik. Berlin: Springer Spektrum, 2018.(Chapter: Beobachtbarkeit der Potenziale)
    \item A. Akhmerov, J. Sau, and B. van Heck, “Tight-binding models in a magnetic field: Peierls substitution,” in Topology in Condensed Matter: Tying Quantum Knots. (URL: https://topocondmat.org/w2-majorana/peierls/) 
\end{itemize}
\end{document}